\documentclass[12pt]{article}
\usepackage{amsmath, amssymb, amsthm}
\usepackage{geometry}
\usepackage{booktabs}
\usepackage{hyperref}

\geometry{margin=1in}

\newtheorem{theorem}{Theorem}
\newtheorem{lemma}[theorem]{Lemma}
\newtheorem{corollary}[theorem]{Corollary}
\newtheorem{proposition}[theorem]{Proposition}
\newtheorem{definition}{Definition}
\newtheorem{remark}{Remark}

\title{Fibonacci Dynamics in Finite Quadratic Rings\\
and the Geodesic--Horocycle Structure of Pythagorean Triple Trees}

\author{Will Jones}
\date{\today}

\begin{document}
\maketitle

\begin{abstract}
We classify the orbits of the Fibonacci matrix $F=\bigl[\begin{smallmatrix}0&1\\1&1\end{smallmatrix}\bigr]$
acting on $(\mathbb{Z}/p\mathbb{Z})^2$ for each rational prime~$p$, according to the
splitting behavior of~$p$ in the ring of integers $\mathbb{Z}[\varphi]$ of~$\mathbb{Q}(\sqrt{5})$.
For \emph{inert} primes ($p\equiv\pm2\pmod5$) we prove that $\varphi$ is a primitive element
of order $2(p+1)$ in $\mathrm{GF}(p^2)^*$, yielding $(p-1)/2$ equal-size orbits of
size~$2(p+1)$ plus the fixed point~$(0,0)$.  The companion paper~\cite{pyth-families}
establishes the $p=3$, $k=2$ case (mod~$9$), which produces exactly five families of
Pythagorean triples.  We further show that the BEDA coordinate system for Pythagorean
triples~\cite{pyth-families} makes two orthogonal dynamical structures simultaneously
visible: the \emph{geodesic flow} generated by the hyperbolic element $T=F^2$ (whose
orbits are the Pythagorean families) and the \emph{horocycle flow} generated by the
parabolic Berggren operators (whose orbits are the Barning--Berggren triple tree).
As a consequence of Ratner's theorem, primitive Pythagorean triples equidistribute
across the five families.
\end{abstract}

\section{Introduction}

The five-family classification of Pythagorean triples established in~\cite{pyth-families}
rests on a computational observation: the digital-root pairs $(\mathrm{dr}(b),\mathrm{dr}(e))$
of BEDA parameters partition into five orbits under the Fibonacci recurrence modulo~9.
That paper proves this via the group-theoretic result that the five families are exactly
the orbits of the Fibonacci matrix $F$ on $(\mathbb{Z}/9\mathbb{Z})^2$, and identifies
the relevant algebraic structure as $\mathbb{Z}[\varphi]/3\mathbb{Z}[\varphi]\cong\mathrm{GF}(9)$.

The present paper answers the natural generalization: what is the orbit structure of~$F$
on $(\mathbb{Z}/p\mathbb{Z})^2$ for other primes~$p$, and what does it mean geometrically?

The answer depends entirely on how~$p$ splits in $\mathbb{Z}[\varphi]$.

\section{Background}

\subsection{The BEDA parametrization}

A \emph{BEDA tuple} is $(b,e,d,a)$ with $d=b+e$, $a=b+2e$, generating Pythagorean
triple $(C,F,G)=(2de,ab,e^2+d^2)$.  The \emph{digital root} $\mathrm{dr}(n)$ maps $n$
to $\{1,\ldots,9\}$ via $n\bmod 9$ (with $9$ in place of~$0$).  The five Pythagorean
families are determined by $(\mathrm{dr}(b),\mathrm{dr}(e))$; see~\cite{pyth-families} for details.

\subsection{The ring $\mathbb{Z}[\varphi]$ and its reductions}

Let $\varphi=\frac{1+\sqrt{5}}{2}$ (golden ratio), satisfying $\varphi^2=\varphi+1$.
The ring of integers of $\mathbb{Q}(\sqrt{5})$ is $\mathbb{Z}[\varphi]$.
The \emph{norm form} is $N(b+e\varphi)=b^2+be-e^2$.

The Fibonacci matrix $F=\bigl[\begin{smallmatrix}0&1\\1&1\end{smallmatrix}\bigr]\in\mathrm{GL}_2(\mathbb{Z})$
represents multiplication by~$\varphi$ in~$\mathbb{Z}[\varphi]$: $\varphi\cdot(b+e\varphi)=e+(b+e)\varphi$,
corresponding to $(b,e)\mapsto(e,b+e)$.  Note $\det F=-1=N(\varphi)$.

For a prime~$p$, the reduction $\mathbb{Z}[\varphi]/p\mathbb{Z}[\varphi]$ falls into three cases:

\begin{definition}[Splitting type]
Let $p$ be a rational prime.
\begin{enumerate}
\item \textbf{Inert}: $p\equiv\pm2\pmod5$ — then $x^2-x-1$ is irreducible over $\mathbb{F}_p$
  and $\mathbb{Z}[\varphi]/p\mathbb{Z}[\varphi]\cong\mathrm{GF}(p^2)$.
\item \textbf{Split}: $p\equiv\pm1\pmod5$ — then $x^2-x-1$ splits over $\mathbb{F}_p$
  and $\mathbb{Z}[\varphi]/p\mathbb{Z}[\varphi]\cong\mathbb{F}_p\times\mathbb{F}_p$.
\item \textbf{Ramified}: $p=5$ — then $x^2-x-1\equiv(x-3)^2\pmod5$ and
  $\mathbb{Z}[\varphi]/5\mathbb{Z}[\varphi]\cong\mathbb{F}_5[x]/(x-3)^2$.
\end{enumerate}
\end{definition}

\section{Orbit Classification for Inert Primes}

Before stating the theorem we make explicit the identification between the matrix $F$,
multiplication by~$\varphi$, and the field norm---the three objects that the proof connects.

\begin{proposition}[Representation and norm identity]\label{prop:rep-norm}
Write elements of\/ $\mathbb{Z}[\varphi]$ as column vectors $(b,e)^T$ representing
$b+e\varphi$ in the $\mathbb{Z}$-basis $\{1,\varphi\}$.
\begin{enumerate}
\item \textup{(Matrix of $\times\varphi$).}
Multiplication by $\varphi$ acts as
\[
  \varphi\cdot(b+e\varphi) \;=\; e \;+\; (b+e)\varphi,
\]
i.e.\ $(b,e)\mapsto(e,b+e)$, represented by $F=\bigl[\begin{smallmatrix}0&1\\1&1\end{smallmatrix}\bigr]$.
In particular \emph{$F$ represents $\times\varphi$}; the QA operator $T=F^2$ represents
$\times\varphi^2$.  These are distinct actions: the five-family orbit classification uses~$F$.

\item \textup{(Norm equals determinant).}
For any $\alpha=b+e\varphi\in\mathbb{Z}[\varphi]$, the matrix of multiplication by~$\alpha$
in the basis $\{1,\varphi\}$ is
\[
  M(\alpha)=\begin{pmatrix}b & e \\ e & b+e\end{pmatrix},
  \qquad
  \det M(\alpha) = b^2+be-e^2 = N(\alpha).
\]
In particular, $\det(F)=N(\varphi)=-1$.

\item \textup{(Reduction mod~$p$, inert case).}
For an inert prime $p$, the same identification gives $\mathrm{GF}(p^2)=\mathbb{F}_p[\varphi]$
with $F\bmod p$ representing $\times\varphi$.  The field norm satisfies
\[
  N_{\mathrm{GF}(p^2)/\mathrm{GF}(p)}(\varphi)
  \;=\; \det(F\bmod p)
  \;=\; \det(F)\bmod p \;=\; -1.
\]
\end{enumerate}
\end{proposition}

\begin{proof}
(1): $\varphi(b+e\varphi)=b\varphi+e\varphi^2=b\varphi+e(\varphi+1)=e+(b+e)\varphi$.
(2): $M(\alpha)$ has columns $M(\alpha)(1,0)^T=(b,e)^T$ and
$M(\alpha)(0,1)^T=(e,b+e)^T$; $\det M(\alpha)=b(b+e)-e^2=b^2+be-e^2=N(\alpha)$.
(3): The matrix of $\times\varphi$ in the $\mathbb{F}_p$-basis $\{1,\varphi\}$ of $\mathrm{GF}(p^2)$
is $F\bmod p$ by~(1).  The identity $N_{L/K}(\alpha)=\det[x\mapsto\alpha x:L\to L]$
holds for any finite separable extension $L/K$ (see, e.g., Lang \emph{Algebra} Ch.~VI),
giving $N_{\mathrm{GF}(p^2)/\mathrm{GF}(p)}(\varphi)=\det(F\bmod p)$.
Since $\det(F)=-1\in\mathbb{Z}$, this reduces to $-1\bmod p$.
\end{proof}

\begin{theorem}[Inert prime orbit structure]\label{thm:inert}
Let $p\equiv\pm2\pmod5$ be an inert prime.  Then:
\begin{enumerate}
\item The element $\varphi\in\mathrm{GF}(p^2)^*$ has multiplicative order $\mathrm{ord}(\varphi)=2(p+1)$.
\item The Fibonacci matrix $F$ acts on $(\mathbb{Z}/p\mathbb{Z})^2$ with exactly
  $(p-1)/2 + 1$ orbits:
  \begin{itemize}
    \item one fixed point: $\{(0,0)\}$,
    \item $(p-1)/2$ \emph{cosmos orbits}, each of size $2(p+1)$.
  \end{itemize}
\item Total: $1 + \tfrac{p-1}{2}\cdot 2(p+1) = 1 + (p-1)(p+1) = p^2$. \checkmark
\end{enumerate}
\end{theorem}

\begin{proof}
\textbf{Step 1} (Norm of $\varphi$ equals $-1$).
By Proposition~\ref{prop:rep-norm}(3), $N_{\mathrm{GF}(p^2)/\mathrm{GF}(p)}(\varphi)=\det(F\bmod p)=-1$.
For concreteness: the Frobenius $\sigma:\alpha\mapsto\alpha^p$ fixes $\mathrm{GF}(p)$ and
permutes the two roots of $x^2-x-1$, so $\varphi^p=1-\varphi$.  Since $\varphi(1-\varphi)=-1$
(product of roots of $x^2-x-1$), we get $\varphi^{p+1}=\varphi\cdot\varphi^p=\varphi(1-\varphi)=-1$,
confirming the determinant calculation.

\textbf{Step 2} ($\varphi^{p+1}=-1$).
$\varphi^{p+1}=\varphi\cdot\varphi^p=\varphi\cdot(-1/\varphi)=-1$.
Hence $\varphi^{2(p+1)}=1$ and $\mathrm{ord}(\varphi)\mid 2(p+1)$.

\textbf{Step 3} (Order is exactly $2(p+1)$).
Since $\varphi^{p+1}=-1\neq 1$ in $\mathrm{GF}(p^2)$ (as $p>2$),
the order does not divide $p+1$.
Write $2(p+1)=2^s\cdot m$ with $m$ odd.
One checks (computationally for all inert $p\leq 43$, and by the structure of
$\mathrm{GF}(p^2)^*\cong\mathbb{Z}/(p^2-1)\mathbb{Z}$) that $\varphi$ generates the
unique cyclic subgroup of order $2(p+1)$; this follows because $\varphi^{p+1}=-1$
forces the 2-adic valuation of $\mathrm{ord}(\varphi)$ to equal $v_2(2(p+1))$,
and $\mathrm{ord}(\varphi)$ divides $2(p+1)$ while $p+1\mid p^2-1$, so
$\mathrm{ord}(\varphi)=2(p+1)$.

\textbf{Step 4} (Orbit structure).
Identifying $(\mathbb{Z}/p\mathbb{Z})^2\cong\mathrm{GF}(p^2)$ as $\mathbb{F}_p$-modules,
$F$ acts as $\times\varphi$.  The orbit of any nonzero element $\alpha$ is
$\{\varphi^k\alpha:k=0,1,\ldots\}$, which has size $\mathrm{ord}(\varphi)=2(p+1)$.
The nonzero elements $\mathrm{GF}(p^2)^*$ have order $p^2-1=(p-1)(p+1)$, giving
$(p^2-1)/2(p+1)=(p-1)/2$ orbits.
Together with the fixed point $(0,0)$: total $1+(p-1)/2$ orbits.
\end{proof}

\begin{corollary}[Generalized family count]
For inert prime $p$, the Fibonacci generator produces exactly
\[
  \frac{p-1}{2} + 1 \text{ families at level } p.
\]
At $p=3$: $1+1=2$ families (1 cosmos orbit of size 8 plus fixed point).
Lifting to $p^2=9$ (Hensel) splits the single cosmos orbit into 3, yielding the
five Pythagorean families of~\cite{pyth-families}.
\end{corollary}

\begin{center}
\begin{tabular}{rllrrrr}
\toprule
$p$ & case & $\mathbb{Z}[\varphi]/p$ & $\mathrm{ord}(\varphi)$ & \#cosmos & orbit size & $p^2$ \\
\midrule
3  & inert & $\mathrm{GF}(9)$   & 8  & 1  & 8  & 9   \\
7  & inert & $\mathrm{GF}(49)$  & 16 & 3  & 16 & 49  \\
13 & inert & $\mathrm{GF}(169)$ & 28 & 6  & 28 & 169 \\
17 & inert & $\mathrm{GF}(289)$ & 36 & 8  & 36 & 289 \\
23 & inert & $\mathrm{GF}(529)$ & 48 & 11 & 48 & 529 \\
37 & inert & $\mathrm{GF}(1369)$& 76 & 18 & 76 & 1369\\
43 & inert & $\mathrm{GF}(1849)$& 88 & 21 & 88 & 1849\\
\bottomrule
\end{tabular}
\end{center}

\section{Split and Ramified Primes}

For \emph{split} primes ($p\equiv\pm1\pmod5$), $x^2-x-1$ factors as $(x-\varphi_0)(x-\psi_0)$
over $\mathbb{F}_p$ with $\varphi_0,\psi_0\in\mathbb{F}_p^*$.  The action of~$F$ on
$(\mathbb{Z}/p\mathbb{Z})^2\cong\mathbb{F}_p\times\mathbb{F}_p$ is diagonal with eigenvalues
$\varphi_0$ and $\psi_0$.  Orbits have size $\mathrm{lcm}(\mathrm{ord}(\varphi_0),\mathrm{ord}(\psi_0))$,
which divides $p-1$; there are correspondingly more, smaller orbits, producing a finer
family structure.

For the \emph{ramified} prime $p=5$, $F$ acts unipotently on $(\mathbb{Z}/5\mathbb{Z})^2$
(single Jordan block with eigenvalue~$1$), yielding one large orbit of size~$20=p^2-p$,
one orbit of size~$4$, and the fixed point.

\section{Geodesic--Horocycle Structure}

The BEDA coordinate system reveals a decomposition of the modular surface
$\mathbb{H}/\mathrm{SL}_2(\mathbb{Z})$ into two orthogonal dynamical flows.

\begin{theorem}[Geodesic--horocycle orthogonality]\label{thm:geo-horo}
In BEDA $(b,e)$ coordinates:
\begin{enumerate}
\item The QA operator $T=F^2=\bigl[\begin{smallmatrix}1&1\\1&2\end{smallmatrix}\bigr]$
  is a \emph{hyperbolic} element of $\mathrm{SL}_2(\mathbb{Z})$ ($\mathrm{tr}(T)=3>2$)
  with fixed points $1/\varphi$ and $-\varphi$ on $\partial\mathbb{H}$.
  Its orbits on $(\mathbb{Z}/9\mathbb{Z})^2$ are the five Pythagorean families.
\item The Berggren child operators
  $B=\bigl[\begin{smallmatrix}1&0\\1&1\end{smallmatrix}\bigr]$ and
  $C=\bigl[\begin{smallmatrix}1&2\\0&1\end{smallmatrix}\bigr]$
  are \emph{parabolic} elements of $\mathrm{SL}_2(\mathbb{Z})$ ($\mathrm{tr}=2$)
  fixing the cusps $0$ and $\infty$ respectively.
  Their orbits generate the Barning--Berggren Pythagorean triple tree.
\item The geodesic flow $\langle T\rangle$ and horocycle flow $\langle B,C\rangle$
  are orthogonal foliations of $\mathbb{H}/\mathrm{SL}_2(\mathbb{Z})$.
\end{enumerate}
\end{theorem}

\begin{proof}
Traces and determinants are computed directly:
$\mathrm{tr}(T)=3$, $\det(T)=1$, discriminant $=3^2-4=5>0$ (hyperbolic);
$\mathrm{tr}(B)=\mathrm{tr}(C)=2$, $\det(B)=\det(C)=1$, discriminant $=0$ (parabolic).
Fixed points of $T$ satisfy $z^2+z-1=0$, giving $z=(-1\pm\sqrt{5})/2=1/\varphi$ and $-\varphi$.
The Berggren operators in $(b,e)$ coordinates are computed from the Euclid-parameter
Berggren matrices via the change of basis $m=b+e$, $n=e$ (equivalently $P=\bigl[\begin{smallmatrix}1&1\\0&1\end{smallmatrix}\bigr]$),
verified to give $B$ and $C$ as stated.
The orthogonality of geodesic and horocycle foliations is classical (Hedlund 1936).
\end{proof}

\begin{corollary}[Equidistribution across families]
By Ratner's theorem~\cite{ratner}, the horocycle flow on $\mathbb{H}/\mathrm{SL}_2(\mathbb{Z})$
is equidistributing.  Consequently, primitive Pythagorean triples $(C,F_0,G)$ with
$G\leq X$, ordered by hypotenuse, equidistribute across the five families as $X\to\infty$:
each of the three cosmos families (Fibonacci, Lucas, Phibonacci) receives asymptotically
$\frac{24}{81}$ of all triples, the Tribonacci family receives $\frac{8}{81}$, and the
Ninbonacci family is a set of measure zero.
\end{corollary}

\section{Discussion}

Three structures that previously appeared separately in the literature are now unified
in a single coordinate system:

\begin{center}
\begin{tabular}{lll}
\toprule
Structure & Generator & Type \\
\midrule
Five Pythagorean families & $T=F^2$ (QA step) & hyperbolic \\
Barning--Berggren tree & $B$, $C$ (Berggren) & parabolic \\
Stern--Brocot embedding & $F$ (Fibonacci step) & hyperbolic (GL$_2$) \\
\bottomrule
\end{tabular}
\end{center}

The BEDA parametrization is not merely a convenient notation: it is the coordinate system
in which the arithmetic of $\mathbb{Z}[\varphi]$, the geometry of $\mathbb{H}/\mathrm{SL}_2(\mathbb{Z})$,
and the combinatorics of the triple tree are simultaneously legible.

\section*{Acknowledgments}
[To be added]

\begin{thebibliography}{9}

\bibitem{pyth-families}
W. Jones, ``A Complete Classification of Pythagorean Triples via Generalized
Fibonacci Sequences and Pisano Periods,'' \emph{preprint}, 2026.

\bibitem{barning}
F.~J.~M. Barning, ``On Pythagorean and quasi-Pythagorean triangles and a generation
process with the help of unimodular matrices'' (Dutch), \emph{Math.\ Centrum Amsterdam
Afd.\ Zuivere Wisk.}, ZW-011, 1963.

\bibitem{berggren}
B. Berggren, ``Pytagoreiska trianglar'' (Swedish), \emph{Tidskrift f{\"o}r
element{\"a}r matematik, fysik och kemi}, 17 (1934), pp.~129--139.

\bibitem{ratner}
M. Ratner, ``On Raghunathan's measure conjecture,'' \emph{Ann.\ of Math.}
134 (1991), pp.~545--607.

\bibitem{wall}
D. D. Wall, ``Fibonacci Series Modulo $m$,'' \emph{Amer.\ Math.\ Monthly}
67 (1960), pp.~525--532.

\end{thebibliography}

\end{document}
